\section{Enforcement and Reconnection}
\label{sec:enforcement}

The authority model says who held command when. This section says what devices do about
it while they cannot reach each other, and what happens to authority and orders when they
reconnect. Three mechanisms matter: the timeout with which a commander's device stops on
its own (\cref{sec:ack}), the moment at which an appointment takes effect
(\cref{sec:acceptance}), and the settlement of orders (\cref{sec:settlement}).

\subsection{Enforcement on Devices}
\label{sec:device-enforcement}

A device lets a person add a decision only if two things hold on the records that device
holds: the person holds command now, and the policy's local condition allows it---the
lease for S1, S2 and S3+L, the commander's timeout $T'$ for S4. The same check guards a
handover, so a commander whose device has stopped cannot hand over either. A decision that
arrives from another device is always kept; if it rests on an authority record the device
has not received, it stays provisional until that record arrives. Field submissions are
always accepted.

\subsection{Acknowledged Timeouts}
\label{sec:ack}

Under S4, B activates once B's device has heard nothing from A's device for $T$, and A's
device stops accepting A's decisions once it has heard nothing from B's for $T' < T$. The
system model argued that A and B ``measure the same silence'', so that $T' < T$ keeps one
device at a time accepting. On links that store and forward, they do not.
\Cref{fig:relay} shows how. A vehicle that lost contact with B still holds a heartbeat
B sent before activating; when A's handheld reaches that vehicle again, the heartbeat
arrives, and A's device, having just ``heard from B'', starts accepting again. The
heartbeat is fresh enough by its timestamp, yet B has meanwhile taken command. In our
witness the two record side by side for \tbRelayOverlap{} (\cref{sec:eval-faults}).

\begin{figure}[t]
\centering
% A stale heartbeat, relayed by V1, revives A's stopped device under S4 (Sec. 5.2).
% Times follow the testbed's relay fault: t0 = cut, t1 = t0 + 200 s, T' = 240 s,
% T = 300 s, Wi-Fi back at t0 + 330 s. One time unit = 50 s.
\begin{tikzpicture}[x=1.02cm, y=0.62cm, every node/.style={font=\scriptsize}]
  \foreach \y/\l in {3/A's handheld, 2/V1 (relay), 1/B (V2)} {
    \draw[black!25] (0,\y) -- (10,\y);
    \node[anchor=east] at (-0.1,\y) {\l};
  }
  \draw[->, black!55] (0,0.05) -- (10.2,0.05) node[right]{time};
  \foreach \x/\l in {0/$t_0$, 4/$t_1$} {
    \draw[black!55] (\x,-0.05) -- (\x,0.15); \node[below] at (\x,0) {\l};
  }
  % A's handheld is cut off from t0 until its Wi-Fi to V1 returns at 6.6.
  \fill[warn!9] (0,2.72) rectangle (6.6,3.28);
  \node[text=warn, font=\scriptsize\bfseries] at (0.15,3) {$\times$};
  % V1 loses B and the core at t1.
  \fill[black!6] (4,1.72) rectangle (10,2.28);
  \node[text=black!50, font=\scriptsize\itshape] at (8.4,2) {V1 cut off from B};
  % B's heartbeats reach V1 until t1.
  \foreach \x in {1,2,3,4} { \draw[->, accept, thin] (\x,1.1) -- (\x,1.88); }
  \node[text=accept, anchor=west] at (4.1,1.5) {B's heartbeats reach V1 until $t_1$};
  % A's device stops at T' after the last heartbeat of B's it received (about t0).
  \draw[warn, line width=1.4pt] (4.8,2.66) -- (4.8,3.34);
  \node[text=warn, above] at (4.8,3.3) {A stops ($T'$)};
  % B activates at T.
  \draw[accept, dashed, thick] (6,0.7) -- (6,1.3);
  \node[text=accept, anchor=north] at (6,0.66) {B activates ($T$)};
  \fill[accept!18] (6,0.72) rectangle (10,1.28);
  \node[text=accept] at (8.6,1) {B accepts};
  % Wi-Fi returns: V1 hands A the heartbeat of B's from t1.
  \draw[->, fence, thick] (6.6,2.12) -- (6.6,2.88);
  \node[text=fence, anchor=east] at (6.52,2.5) {V1 relays B's beat from $t_1$};
  % A's device accepts again until t1 + T' (8.8).
  \fill[fence!25] (6.6,2.72) rectangle (8.8,3.28);
  \node[text=fence] at (7.7,3) {A accepts again};
  \draw[fence, line width=1.4pt] (8.8,2.66) -- (8.8,3.34);
  \node[text=fence, above] at (8.8,3.3) {$t_1 + T'$};
\end{tikzpicture}

\caption{A stale heartbeat revives a stopped commander. A's handheld is cut off at $t_0$.
V1 last hears B at $t_1$ and keeps B's heartbeat from then. A's device stops at $T'$ after
the last beat of B's it received, and B activates at $T$. When A's Wi-Fi to V1 returns, V1
hands A B's heartbeat from $t_1$, younger than $T'$; A's device accepts again, beside B,
until $t_1 + T'$. With an acknowledged timeout, the same heartbeat reports that B last
heard A at $t_0$, and A's device stays stopped.}
\label{fig:relay}
\Description{A timeline with three lanes: A's handheld, vehicle V1 and B's vehicle. A is cut off at t0 and stops at T-prime; B activates at T; V1, which last heard B at t1, later hands A that old heartbeat, and A accepts again beside B until t1 plus T-prime.}
\end{figure}

The fix is the one leases have long used against this: renew by acknowledgment, not by
hearing. Every heartbeat carries, for each other device, the timestamp of the latest
heartbeat the sender has received from it. A's device stops once the latest of \emph{its
own} heartbeats that B's device is known to have received is more than $T'$ old.

\begin{proposition}[One device at a time]
\label{prop:single-acceptor}
Let clock offsets between devices differ by at most $\epsilon$, and let a device that
holds a heartbeat hold every record its sender held when sending it. Let B's device
activate at the first time $t_B$ at which $C_B(t_B) - y \geq T$, where $y$ is the timestamp
of the latest heartbeat of A's it holds, and let A's device accept A's decisions at time
$t$ only if A holds command on its records and $C_A(t) - x(t) \leq T'$, where $x(t)$ is the
largest timestamp of A's heartbeats that the heartbeats of B's held by A's device report B
to have received. If $T - T' > \epsilon$, A's device accepts none of A's decisions at or
after $t_B$.
\end{proposition}

\begin{proof}[Proof sketch]
If A's device holds a heartbeat of B's sent at or after $t_B$, it holds B's activation, so A
no longer holds command on its records and the device refuses. Otherwise every heartbeat
of B's that it holds was sent before $t_B$, and reports a value at most $y_B$, the value of
$y$ at $t_B$; hence $x(t) \leq y_B$. Accepting at $t$ needs $C_A(t) \leq y_B + T'$, while
activation gives $C_B(t_B) \geq y_B + T$. With $C_v(t) = t + o_v$ and
$o_B - o_A \leq \epsilon$, $t \leq t_B - (T - T') + \epsilon < t_B$.
\end{proof}

The margin needs to cover only the clock error, not message delay, the heartbeat period or
the path through the mesh: those make $x$ lag behind $y$, which only stops A earlier. The
rule the system model used, with $x$ replaced by the timestamp of the latest heartbeat of
B's, has no such bound: a heartbeat of B's sent after B last heard A, but before B
activated, can reach A arbitrarily late, as in \cref{fig:relay}. Drift adds a term
proportional to $T$, which is negligible while GNSS is available. Leases renewed at the
regional server (S1, S2, S3+L) have the same weakness and the same fix.
\pending{Model-check \cref{prop:single-acceptor} in TLA+ with relaying channels.}

\subsection{When an Appointment Takes Effect}
\label{sec:acceptance}

A takeover by C, or a promotion by an operator, is signed at the regional centre. The
system model let it take effect when signed. When the scene is cut off from the centre, as
in P-B, this rewrites the past (\cref{fig:acceptance}a). A is fit and does not answer the
call, so C takes over; A, unaware, keeps giving orders over the mesh, and crews see them as
valid. When cellular returns, the takeover arrives with its earlier timestamp, and every
order A gave since becomes residual---orders that A gave in good faith and crews have
carried out. In the evaluation this turns \tbSignedResidual{} orders residual, every one of
which the three crew devices had shown as valid: \tbSignedFlips{} flips
(\cref{sec:eval-faults}).

\begin{figure}[t]
\centering
% A takeover signed at the centre while the scene is cut off (P-B), effective when signed
% (a) or when the appointee's device accepts it (b). Sec. 5.3. Colours from main.tex.
\newcommand{\acclanes}{%
  \foreach \y/\l in {3/A (mesh to crews), 2/B (V2), 1/C (centre)} {
    \draw[black!25] (0,\y) -- (10,\y);
    \node[anchor=east] at (-0.1,\y) {\l};
  }
  \fill[warn!8] (1,0.72) rectangle (8,3.28);
  \node[text=warn, font=\scriptsize\itshape, anchor=north west] at (1.05,0.68)
    {P-B: cellular down, mesh up};
  \fill[accept] (3,1) circle (2.2pt);
  \node[text=accept, anchor=west] at (3.1,1.4) {C takes over, appointing B ($s$)};
  \draw[->, black!60, thick] (8,1.15) -- (8,2.85);
  \node[anchor=north] at (8,0.68) {heal ($h$)};
}
\begin{tikzpicture}[x=0.85cm, y=0.58cm, every node/.style={font=\scriptsize}]
  \node[anchor=west, font=\footnotesize\bfseries] at (-1.6,3.95) {(a) effective when signed};
  \acclanes
  \fill[resid] (4.6,3) circle (2.6pt); \node[above] at (4.6,3.1) {d1};
  \fill[resid] (6.4,3) circle (2.6pt); \node[above] at (6.4,3.1) {d2};
  \fill[accept!18] (3,1.72) rectangle (10,2.28);
  \node[text=accept] at (9.1,2) {B holds};
  \node[text=resid, anchor=west, align=left] at (8.15,3.05) {d1, d2 turn\\residual};
\end{tikzpicture}

\vspace{6pt}
\begin{tikzpicture}[x=0.85cm, y=0.58cm, every node/.style={font=\scriptsize}]
  \node[anchor=west, font=\footnotesize\bfseries] at (-1.6,3.95) {(b) effective when B's device accepts};
  \acclanes
  \fill[good] (4.6,3) circle (2.6pt); \node[above] at (4.6,3.1) {d1};
  \fill[good] (6.4,3) circle (2.6pt); \node[above] at (6.4,3.1) {d2};
  \fill[accept!18] (8,1.72) rectangle (10,2.28);
  \fill[accept] (8,2) circle (2.2pt);
  \node[text=accept] at (9.1,2) {B holds};
  \node[text=good, anchor=west, align=left] at (8.15,3.05) {d1, d2 stand;\\A learns at $h$};
\end{tikzpicture}

\caption{A takeover that cannot reach the scene. In P-B the scene keeps its mesh but loses
the regional centre. C, unable to reach A by voice, signs a takeover appointing B at $s$.
(a) Effective when signed: A's orders d1 and d2, given over the mesh and seen as valid by
crews, become residual when the links return. (b) Effective when B's device accepts it:
command passes at the heal $h$, d1 and d2 stand, and A's device learns of the change at
once.}
\label{fig:acceptance}
\Description{Two timelines with lanes for A, B and C during a P-B partition. C takes over at s. In panel a, A's orders d1 and d2 turn residual when links heal; in panel b, B accepts the takeover at the heal and d1 and d2 stand.}
\end{figure}

We let a takeover or promotion take effect when the appointee's device accepts it, as a
handover already does when the recipient accepts it. This is the confirmation step of an
interlocked transfer of control~\cite{flemisch2012balance,gil2019designing}, made explicit
for the case where the confirmation itself is delayed. Three rules come with it: the
appointee's device accepts as soon as it holds the appointment and the appointee is fit; an
appointment older than the one in force can no longer be accepted, so a stale takeover
cannot override a newer one; and settlement waits for anyone named in an appointment not
yet accepted (\cref{sec:settlement}). Without the third rule, our search found confirmed
verdicts that a late acceptance later revoked. In P-A, acceptance changes nothing: B
accepts at once, and A, who cannot be told, keeps recording until it reconnects. Those
residual orders are the irreducible cost of a commander who cannot be reached.
\pending{Decision for the supervisors: effective on acceptance.}

\subsection{Settlement}
\label{sec:settlement}

A valid order is provisional until the regional server confirms it; devices that receive
the confirmation, and hold the records it rests on, mark it authoritative. The server
decides nothing: it confirms only that no record that could change the verdict can still
arrive. Per-person sequence numbers make this checkable: a heartbeat from a person's
device carries that person's latest sequence number, so once the server holds a heartbeat
with timestamp at least $\tau + \epsilon$ and every record of that person up to its sequence
number, it knows that person signed nothing it lacks up to $\tau$.

Whose records must be complete? The system model's rule, read literally, names everyone
named in any authority record. Then a commander whose device is lost blocks the settlement
of every later decision, forever (\cref{sec:eval-settlement}). But an unconditional
act---a grant, a takeover or promotion once in effect---overrides everything before it, so
only some people can still change a verdict.

\begin{proposition}[Settlement]
\label{prop:settlement}
For a decision $d$ with timestamp $\tau_d$, let $R(d)$ hold the appointing officer and the
operators, every holder of command since the latest unconditional act before
$\tau_d - \epsilon$, the deputies named in rules for those holders, and anyone named in an
appointment not yet accepted. Let the server confirm $d$ once its records classify $d$ as
valid and, for every person in $R(d)$, it holds that person's records complete up to
$\tau_d + \epsilon$. Under the assumptions of \cref{sec:system-model}:
(i) no device that has marked $d$ authoritative later classifies it otherwise; and
(ii) if $d$ is valid and the devices of everyone in $R(d)$ eventually reach the server, $d$
is confirmed.
\end{proposition}

\begin{proof}[Proof sketch]
The verdict on $d$ depends only on authority acts with timestamps up to $\tau_d + \epsilon$:
those that decide who held command at $\tau_d$, and those within $\epsilon$ of it. Any act
that precedes $d$ in real time has a timestamp below $\tau_d + \epsilon$. An act that changes
who held command after the latest unconditional act can be signed only by the appointing
officer, the holder (a handover), a deputy of the holder (an activation), or an appointee
(an acceptance), all of whom are in $R(d)$; acts before the unconditional act are overridden
by it. Completeness of their records means no such act can still arrive, so the server's
verdict is final, and a device that holds the confirmation and the records it lists computes
the same verdict. For (ii), every completeness condition is eventually met once the devices
of $R(d)$ reach the server.
\end{proof}

Part (ii) is where policies differ. Under S1 and S3 the successor holds command by an
unconditional act, so the predecessor leaves $R(d)$ and the successor's orders settle
within seconds, even if the predecessor's device is never seen again. Under S4 the
successor holds command by an activation, which is valid only if the predecessor was still
the holder, and only the predecessor's records can show that. The deputy's orders therefore
stay provisional until the predecessor's device reconnects, or for ever if it is lost,
unless C ratifies the activation with a takeover. \pending{Model-check part (i) in TLA+;
decide with the supervisors whether S4 requires ratification by C.}

\subsection{What Commanders and Crews See}
\label{sec:readers}

R4 asks that every person who acts on an order can see whether it stands. Each device shows,
for every order, who issued it, under which act, and its status. Orders move between
statuses as records arrive: provisional to authoritative when settled; provisional to
residual when a record arrives showing that its issuer no longer held command; provisional
to in doubt when a change of command turns out to lie within $\epsilon$. A move from
provisional to residual is a flip: a crew may already have acted. Flips cannot be prevented
while a commander cannot be told, but policies differ widely in how many they produce, and
appointments that take effect on acceptance remove those that a takeover from afar
would cause. The evaluation counts flips on the crews' devices.
\pending{Design of the status display, and whether officers read it as intended: expert
walkthrough, \cref{sec:eval-walkthrough}.}
